
\documentclass{article}
\usepackage{colortbl}
\usepackage{makecell}
\usepackage{multirow}
\usepackage{supertabular}

\begin{document}

\newcounter{utterance}

\centering \large Interaction Transcript for game `wordle', experiment `medium\_frequency\_words\_no\_clue\_no\_critic', episode 11 with reflexion{-}Qwen3.5{-}27B.
\vspace{24pt}

{ \footnotesize  \setcounter{utterance}{1}
\setlength{\tabcolsep}{0pt}
\begin{supertabular}{c@{$\;$}|p{.15\linewidth}@{}p{.15\linewidth}p{.15\linewidth}p{.15\linewidth}p{.15\linewidth}p{.15\linewidth}}
    \# & \multicolumn{2}{c}{Player} && \multicolumn{2}{c}{Game Master} \\
    \hline

    \theutterance \stepcounter{utterance}  
    & & & \multicolumn{4}{p{0.6\linewidth}}{
        \cellcolor[rgb]{0.9,0.9,0.9}{
            \makecell[{{p{\linewidth}}}]{
                \texttt{\tiny{[P1$\langle$GM]}}
                \texttt{You have attempted this same task before and did not succeed. These are the plans you wrote after those attempts.} \\
\texttt{Trial \#1: New plan: The previous attempt failed because I included my internal reasoning process (e.g., "wait, 'a' was yellow...", "Let's try...") directly in the "explanation" field, causing the response to become disorganized and likely violating the "concise explanation" rule or the strict format. Additionally, I did not provide a final "guess:" line before the text cut off. Next time, I will strictly separate my internal logic from the output. I will perform the deduction mentally to identify the correct word (likely "taint" or "tania" based on the constraints: 'a' not in 1,2,3; 't' not in 4; 'i' not in 2; 'n' not in 3), and then output *only* a clean, concise explanation of the final choice followed immediately by the "guess:" tag and the five{-}letter word. I will ensure the explanation summarizes the logic without showing the step{-}by{-}step trial{-}and{-}error or self{-}correction.} \\
\\ 
\texttt{Trial \#2: New plan: The previous attempt failed because I included my internal trial{-}and{-}error reasoning (e.g., "Wait...", "Let's try...", "No...") within the "explanation" field, violating the requirement for a concise explanation, and I failed to output the final "guess:" line before the response was cut off. Next time, I will perform all logical deductions and word validation internally before generating the response. I will ensure the "explanation" field contains only a clear, final summary of why the chosen word fits the constraints (e.g., "Based on feedback, 't', 'a', 'i', 'n' are present but misplaced; 'taint' places them in valid positions while avoiding excluded spots"). I will strictly adhere to the format by providing the explanation followed immediately by a valid five{-}letter word on the "guess:" line, ensuring no internal monologue or incomplete sentences appear in the output.} \\
\\ 
\texttt{Trial \#3: New plan: The previous attempt failed because the "explanation" field contained unstructured internal monologue, self{-}correction, and incomplete sentences, violating the requirement for a concise explanation, and the response was cut off before providing the mandatory "guess:" line. Next time, I will perform all logical deductions and word validation internally to determine the correct word (e.g., "taint" or "tania" based on the constraints: t, a, i, n present but not in positions 1, 3, 4, 5 respectively). I will then output *only* a clean, single{-}paragraph explanation summarizing the final logic, followed immediately by the "guess:" tag and a valid five{-}letter word, ensuring no internal reasoning steps or incomplete thoughts appear in the final output.} \\
\\ 
\\ 
\texttt{You are a language wizard who likes to guess words by using the given rules.} \\
\\ 
\texttt{Welcome to Wordle! You have six attempts to guess the target word, a valid English word of five lowercase letters (a{-}z). Please use the tags "explanation:" and "guess:" to provide a concise explanation for each guess.} \\
\\ 
\texttt{For instance, if your guess is "apple", your response should be} \\
\texttt{explanation: this is a common five{-}letter English word, and I am starting my guess with this word.} \\
\texttt{guess: apple} \\
\\ 
\texttt{After each guess, your answer will be validated, and you will receive feedback indicating which letters are correct (green), which letters are correct but in the wrong position (yellow), and which letters are incorrect (red). This feedback can be useful in determining which letters to include or exclude in your next guess.} \\
\\ 
\texttt{For example, the feedback for "apple" might be:} \\
\texttt{guess\_feedback: a<yellow> p<yellow> p<green> l<yellow> e<red>} \\
\\ 
\texttt{The explanation should contain details about how the guess\_feedback is used to arrive at a new guess.} \\
\\ 
\texttt{Do not generate any other text.} \\
\\ 
\texttt{Let's begin with your first guess.} \\
            }
        }
    }
     \\ \\

    \theutterance \stepcounter{utterance}  
    & \multicolumn{4}{p{0.6\linewidth}}{
        \cellcolor[rgb]{0.561,0.737,0.561}{
            \makecell[{{p{\linewidth}}}]{
                \texttt{\tiny{[P1$\rangle$GM]}}
                \texttt{explanation: I am starting with a common five{-}letter word containing frequent vowels and consonants to maximize the information gained from the initial feedback.} \\
\texttt{guess: slate} \\
            }
        }
    }
    & & \\ \\

    \theutterance \stepcounter{utterance}  
    & & & \multicolumn{2}{p{0.3\linewidth}}{
        \cellcolor[rgb]{0.9,0.9,0.9}{
            \makecell[{{p{\linewidth}}}]{
                \texttt{\tiny{[GM$|$GM]}}
                \texttt{attempts = 1} \\
\texttt{target\_word = faint} \\
\texttt{guess = slate} \\
\texttt{guess\_feedback = s<red> l<red> a<yellow> t<yellow> e<red>} \\
            }
        }
    }
    & & \\ \\

    \theutterance \stepcounter{utterance}  
    & & & \multicolumn{4}{p{0.6\linewidth}}{
        \cellcolor[rgb]{0.9,0.9,0.9}{
            \makecell[{{p{\linewidth}}}]{
                \texttt{\tiny{[P1$\langle$GM]}}
                \texttt{guess\_feedback: s<red> l<red> a<yellow> t<yellow> e<red>} \\
\\ 
\texttt{Provide your response only in this format.} \\
\texttt{explanation: details} \\
\texttt{guess: word} \\
            }
        }
    }
     \\ \\

    \theutterance \stepcounter{utterance}  
    & \multicolumn{4}{p{0.6\linewidth}}{
        \cellcolor[rgb]{0.561,0.737,0.561}{
            \makecell[{{p{\linewidth}}}]{
                \texttt{\tiny{[P1$\rangle$GM]}}
                \texttt{explanation: The feedback indicates 'a' and 't' are in the word but not in positions 3 and 4 respectively, while 's', 'l', and 'e' are not in the word. I will choose a word that places 'a' and 't' in new positions and uses common remaining letters like 'r', 'i', and 'n' to test further.} \\
\texttt{guess: train} \\
            }
        }
    }
    & & \\ \\

    \theutterance \stepcounter{utterance}  
    & & & \multicolumn{2}{p{0.3\linewidth}}{
        \cellcolor[rgb]{0.9,0.9,0.9}{
            \makecell[{{p{\linewidth}}}]{
                \texttt{\tiny{[GM$|$GM]}}
                \texttt{attempts = 2} \\
\texttt{target\_word = faint} \\
\texttt{guess = train} \\
\texttt{guess\_feedback = t<yellow> r<red> a<yellow> i<yellow> n<yellow>} \\
            }
        }
    }
    & & \\ \\

    \theutterance \stepcounter{utterance}  
    & & & \multicolumn{4}{p{0.6\linewidth}}{
        \cellcolor[rgb]{0.9,0.9,0.9}{
            \makecell[{{p{\linewidth}}}]{
                \texttt{\tiny{[P1$\langle$GM]}}
                \texttt{guess\_feedback: t<yellow> r<red> a<yellow> i<yellow> n<yellow>} \\
\\ 
\texttt{Provide your response only in this format.} \\
\texttt{explanation: details} \\
\texttt{guess: word} \\
            }
        }
    }
     \\ \\

    \theutterance \stepcounter{utterance}  
    & \multicolumn{4}{p{0.6\linewidth}}{
        \cellcolor[rgb]{0.561,0.737,0.561}{
            \makecell[{{p{\linewidth}}}]{
                \texttt{\tiny{[P1$\rangle$GM]}}
                \texttt{explanation: The feedback confirms 't', 'a', 'i', and 'n' are in the word but misplaced in "train" (positions 1, 3, 2, 5), while 'r' is excluded. I need to rearrange these four letters into valid positions and add a new consonant. Placing 't' at the end, 'a' at the start, 'i' in the middle, and 'n' at the end creates "taint", which fits the constraints of avoiding the previous positions for these letters.} \\
\texttt{guess: taint} \\
            }
        }
    }
    & & \\ \\

    \theutterance \stepcounter{utterance}  
    & & & \multicolumn{2}{p{0.3\linewidth}}{
        \cellcolor[rgb]{0.9,0.9,0.9}{
            \makecell[{{p{\linewidth}}}]{
                \texttt{\tiny{[GM$|$GM]}}
                \texttt{attempts = 3} \\
\texttt{target\_word = faint} \\
\texttt{guess = taint} \\
\texttt{guess\_feedback = t<red> a<green> i<green> n<green> t<green>} \\
            }
        }
    }
    & & \\ \\

    \theutterance \stepcounter{utterance}  
    & & & \multicolumn{4}{p{0.6\linewidth}}{
        \cellcolor[rgb]{0.9,0.9,0.9}{
            \makecell[{{p{\linewidth}}}]{
                \texttt{\tiny{[P1$\langle$GM]}}
                \texttt{guess\_feedback: t<red> a<green> i<green> n<green> t<green>} \\
\\ 
\texttt{Provide your response only in this format.} \\
\texttt{explanation: details} \\
\texttt{guess: word} \\
            }
        }
    }
     \\ \\

    \theutterance \stepcounter{utterance}  
    & \multicolumn{4}{p{0.6\linewidth}}{
        \cellcolor[rgb]{0.561,0.737,0.561}{
            \makecell[{{p{\linewidth}}}]{
                \texttt{\tiny{[P1$\rangle$GM]}}
                \texttt{explanation: The feedback shows 'a', 'i', 'n', and the second 't' are in the correct positions (2, 3, 4, 5), but the first 't' is incorrect. This means the word ends in "aint" and the first letter is not 't'. I will try a common word starting with a different letter that fits the pattern "\_aint", such as "paint".} \\
\texttt{guess: paint} \\
            }
        }
    }
    & & \\ \\

    \theutterance \stepcounter{utterance}  
    & & & \multicolumn{2}{p{0.3\linewidth}}{
        \cellcolor[rgb]{0.9,0.9,0.9}{
            \makecell[{{p{\linewidth}}}]{
                \texttt{\tiny{[GM$|$GM]}}
                \texttt{attempts = 4} \\
\texttt{target\_word = faint} \\
\texttt{guess = paint} \\
\texttt{guess\_feedback = p<red> a<green> i<green> n<green> t<green>} \\
            }
        }
    }
    & & \\ \\

    \theutterance \stepcounter{utterance}  
    & & & \multicolumn{4}{p{0.6\linewidth}}{
        \cellcolor[rgb]{0.9,0.9,0.9}{
            \makecell[{{p{\linewidth}}}]{
                \texttt{\tiny{[P1$\langle$GM]}}
                \texttt{guess\_feedback: p<red> a<green> i<green> n<green> t<green>} \\
\\ 
\texttt{Provide your response only in this format.} \\
\texttt{explanation: details} \\
\texttt{guess: word} \\
            }
        }
    }
     \\ \\

    \theutterance \stepcounter{utterance}  
    & \multicolumn{4}{p{0.6\linewidth}}{
        \cellcolor[rgb]{0.561,0.737,0.561}{
            \makecell[{{p{\linewidth}}}]{
                \texttt{\tiny{[P1$\rangle$GM]}}
                \texttt{explanation: The feedback confirms the pattern is "\_aint" with 'a', 'i', 'n', 't' in positions 2, 3, 4, 5. The first letter is not 't' (from "taint") and not 'p' (from "paint"). I will try "quaint" but it is six letters, so I need a five{-}letter word. Common words fitting "\_aint" include "stain", "drain", "brain", "grain", "plain", "chain", "frain" (not a word), "swain". Since 's' was red in the first guess "slate", "stain" is out. 'r' was red in "train", so "drain", "brain", "grain" are out. 'p' was red in "paint", so "plain" is out. 'c' and 'h' are untested. "Chain" fits the pattern and uses untested letters.} \\
\texttt{guess: chain} \\
            }
        }
    }
    & & \\ \\

    \theutterance \stepcounter{utterance}  
    & & & \multicolumn{2}{p{0.3\linewidth}}{
        \cellcolor[rgb]{0.9,0.9,0.9}{
            \makecell[{{p{\linewidth}}}]{
                \texttt{\tiny{[GM$|$GM]}}
                \texttt{attempts = 5} \\
\texttt{target\_word = faint} \\
\texttt{guess = chain} \\
\texttt{guess\_feedback = c<red> h<red> a<yellow> i<yellow> n<yellow>} \\
            }
        }
    }
    & & \\ \\

    \theutterance \stepcounter{utterance}  
    & & & \multicolumn{4}{p{0.6\linewidth}}{
        \cellcolor[rgb]{0.9,0.9,0.9}{
            \makecell[{{p{\linewidth}}}]{
                \texttt{\tiny{[P1$\langle$GM]}}
                \texttt{guess\_feedback: c<red> h<red> a<yellow> i<yellow> n<yellow>} \\
\\ 
\texttt{Provide your response only in this format.} \\
\texttt{explanation: details} \\
\texttt{guess: word} \\
            }
        }
    }
     \\ \\

    \theutterance \stepcounter{utterance}  
    & \multicolumn{4}{p{0.6\linewidth}}{
        \cellcolor[rgb]{0.561,0.737,0.561}{
            \makecell[{{p{\linewidth}}}]{
                \texttt{\tiny{[P1$\rangle$GM]}}
                \texttt{explanation: The feedback for "chain" shows 'c' and 'h' are incorrect, and 'a', 'i', 'n' are yellow, meaning they are in the word but not in positions 2, 3, 4. However, previous feedback from "paint" confirmed 'a', 'i', 'n', 't' are green in positions 2, 3, 4, 5. This implies the target word is "\_aint". The first letter cannot be 't', 'p', 'c', 'h', 's', 'r', 'l', 'e'. I need a word starting with a valid letter. "Quaint" is 6 letters. "Saint" has 's' (red). "Drain" has 'r' (red). "Brain" has 'r' (red). "Grain" has 'r' (red). "Swain" has 's' (red). "Waint" is not a word. "Laint" is not a word. "Moint" no. "Vaint" no. Wait, let's re{-}evaluate "chain" feedback. If "paint" gave a<green> i<green> n<green> t<green>, then the word IS \_aint. "chain" gave a<yellow> i<yellow> n<yellow>. This is a contradiction unless the target word is NOT \_aint. Let's re{-}read "paint" feedback: p<red> a<green> i<green> n<green> t<green>. This means the word is definitely \_aint. Why did "chain" give yellow for a, i, n? Because in "chain", a is at 2, i at 3, n at 4. If the target is \_aint, a is at 2, i at 3, n at 4. They should be green. The only explanation is that the target word is NOT \_aint, or I misinterpreted the previous feedback. Let's re{-}read "taint" feedback: t<red> a<green> i<green> n<green> t<green>. This means positions 2,3,4,5 are a,i,n,t. So the word is ?aint. "paint" confirmed p is wrong. "chain" feedback says a,i,n are yellow. This implies in the target word, a is NOT at 2, i is NOT at 3, n is NOT at 4. This contradicts "taint" and "paint" feedback where they were green. Is it possible the target word has multiple instances? No, "taint" has one 'a', one 'i', one 'n'. If they were green, they are in those spots. If "chain" makes them yellow, it means the target word does NOT have a at 2, i} \\
            }
        }
    }
    & & \\ \\

    \theutterance \stepcounter{utterance}  
    & & & \multicolumn{2}{p{0.3\linewidth}}{
        \cellcolor[rgb]{0.9,0.9,0.9}{
            \makecell[{{p{\linewidth}}}]{
                \texttt{\tiny{[GM$|$GM]}}
                \texttt{Error: The response should contain the 'guess:' keyword exactly once.} \\
            }
        }
    }
    & & \\ \\

    \theutterance \stepcounter{utterance}  
    & & & \multicolumn{2}{p{0.3\linewidth}}{
        \cellcolor[rgb]{0.9,0.9,0.9}{
            \makecell[{{p{\linewidth}}}]{
                \texttt{\tiny{[GM$|$GM]}}
                \texttt{game\_result = ABORT} \\
            }
        }
    }
    & & \\ \\

    \theutterance \stepcounter{utterance}  
    & & \multicolumn{4}{p{0.6\linewidth}}{
        \cellcolor[rgb]{0.95,0.95,0.95}{
            \makecell[{{p{\linewidth}}}]{
                \texttt{\tiny{[P1$\leftrightarrow$P]}}
                \texttt{Post{-}episode Reflexion (not sent to GM):} \\
\texttt{New plan: The previous attempt failed because the "explanation" field devolved into unstructured internal monologue, self{-}correction, and confusion about contradictory feedback, violating the requirement for a concise explanation. Furthermore, the response was cut off before providing the mandatory "guess:" line, causing an ABORT. Next time, I will perform all logical deductions, resolve any apparent contradictions in feedback mentally, and validate the candidate word internally before generating the response. I will ensure the "explanation" field contains only a single, clear sentence summarizing the final logic (e.g., "Based on the feedback, the word must be 'taint' as it places the confirmed letters in their correct positions while avoiding excluded letters"). I will strictly adhere to the format by immediately following the explanation with the "guess:" tag and a valid five{-}letter word, ensuring no internal reasoning, incomplete sentences, or missing tags appear in the output.} \\
            }
        }
    }
    & \\ \\

\end{supertabular}
}

\end{document}
